\chapter{Implementation and Validation}
\label{chapter3}

% ──────────────────────────────────────────────────────────────────────────────────────────────────────────
\section{L-System Interpreter}
\label{sec:interpreter}

\subsection{Symbol Trait}

The interpreter is built around two abstractions: the \texttt{Symbol} trait and the \texttt{Rule} enum. Any
type that implements \texttt{Symbol} may be used as the alphabet of an L-System.

\begin{verbatim}
pub trait Symbol: PartialEq + Copy + Clone {
    fn symbol_type(&self) -> SymbolType; // Terminal | NonTerminal

    fn context(&self) -> ContextAction { // Ignore | Consider | BranchStart | BranchEnd
        ContextAction::Consider
    }
}
\end{verbatim}

Only \texttt{NonTerminal} symbols are subject to rewriting; \texttt{Terminal} symbols pass through unchanged.
This distinction avoids the need for explicit identity productions and is consistent with the standard
L-System formulation. The \texttt{context()} method is used to determine the matching behaviour for
context-sensitive rules, by default all symbols are considered.

\subsection{Rule Variants}

The \texttt{Rule} enum provides four variants covering all required grammar classes:

\begin{verbatim}
pub enum Rule<'a, S: Symbol> {
    Normal(S, &'a [S]),
    Stochastic(S, f32, &'a [S]),
    ContextSensitive(S, Option<S>, Option<S>, &'a [S]),
    Parametric(S, &'a dyn Fn(&S, &mut Vec<S>) -> bool),
}
\end{verbatim}

The type parameter \texttt{<'a, S: Symbol>} makes each rule generic over the grammar alphabet. The Rust trait
bound \texttt{S: Symbol} enforces at compile time that only types implementing \texttt{Symbol} may serve as
alphabet elements. The lifetime \texttt{'a} ties the borrows to the rule~\cite{TraitLifetimeBounds}: each
\texttt{\&'a [S]} is a reference to a caller-allocated array rather than an owned \texttt{Vec<S>},
eliminating the need for heap allocation. This is viable because the output of a \texttt{Normal},
\texttt{Stochastic} or \texttt{ContextSensitive} rule is fixed at grammar-construction time. The
\texttt{Parametric} variant differs as its output is dependent on the parameters passed: its closure writes
into a \texttt{\&mut Vec<S>} provided by the rewriting engine at runtime, avoiding a per-symbol allocation
while still permitting computed output.

\medskip
\texttt{Normal} rules can be used for deterministic, context-free productions, they have two parameters; the
symbol they match and the set of symbols that it is replaced with. The production $A \rightarrow AB$ would be
modelled as \texttt{Rule::Normal(A, \&[A, B])}.

\medskip
\texttt{Stochastic} rules are similar to \texttt{Normal} rules but take an additional probability parameter
representing the likelihood for that production to be applied. A production with a 50\% chance of being
applied, $A[0.5] \rightarrow AB$, is denoted as \texttt{Rule::Stochastic(A, 0.5, \&[A, B])}.

\medskip
\texttt{ContextSensitive} rules take two additional parameters, the left and right context, the rule will
only be applied when the matched symbol appears in between the left and right context. Both the left and
right parameters are optional allowing this rule to model 1L-Systems and 2L-Systems. The production $B < A >
C \rightarrow BC$ is written as \texttt{Rule::ContextSensitive(A, Some(B), Some(C), \&[B, C])}. Left context
is primarily used to model acropetal signals, signals which travel from the base of the plant up, because of
this if a branch is encountered whilst checking for left context it will be skipped. Right context models
basipetal signals, those which travel down the plant, and as a result must check every branch after the symbol.

\medskip
\texttt{Parametric} rules carry a trait object reference (\texttt{\&'a dyn}) for a closure (\texttt{Fn(..) ->
bool}) which takes a shared reference to the matched symbol (\texttt{\&S}) and a mutable reference to a
vector (\texttt{\&mut Vec<S>}) as its parameters and returns a boolean indicating whether it matched the
symbol. As rules can be created dynamically the closure they reference cannot be known until runtime. Trait
objects allow for the late binding of methods where, as described in the Rust Reference~\cite{TraitObjectTypes}:

\begin{displayquote}
  Calling a method on a trait object results in virtual dispatch at runtime: that is, a function pointer is
  loaded from the trait object vtable and invoked indirectly
\end{displayquote}

The closure receives a shared reference to the matched symbol (allowing it to read inline parameters) and
appends its production into an existing \texttt{Vec}. For example the production $I(l): l < 10 \rightarrow
I(l * k)$, which models an internode \texttt{I(l)} that grows by factor~$k$ until it reaches $10$ is written:

\begin{verbatim}
Rule::Parametric(I(0.0), &|s, out| {
    if let &I(l) = s && l < 10 {
        out.push(I(l * k));
        true
    } else {
        false
    }
})
\end{verbatim}

When symbols are matched to rules \texttt{std::mem::discriminant} is used, this means that the rule head
\texttt{I(0.0)} will match any \texttt{I(..)} regardless of the parameters value.~\footnote{Matching
  behaviour technically depends on the implementation of the \texttt{PartialEq} trait for an object
  implementing \texttt{Symbol} which may vary depending on the underlying type, however
\texttt{std::mem::discriminant} is used by all grammars in this project}

\subsection{L-System}

From these abstractions the \texttt{LSystem} struct is defined as:

\begin{verbatim}
pub struct LSystem<'a, S: Symbol> {
    state: Vec<S>,
    rules: Vec<Rule<'a, S>>,
}
\end{verbatim}

\texttt{LSystem::new(axiom, rules)} takes the axiom as a \texttt{\&[S]} slice and the rules as a
\texttt{Vec<Rule>}. \texttt{step()} performs one rewriting pass: each symbol is matched against the rule
list, with context-sensitive rules taking priority, the first matching rule appends its production to the
output; if no rule matches, the symbol is copied unchanged. \texttt{evolve(n)} calls \texttt{step} $n$ times
and \texttt{current()} returns an immutable slice of the current string.

\subsection{Seeded RNG}

Stochastic rules use a thread-local \texttt{SmallRng} from the \texttt{rand} crate, seeded before each scene
build from the scene's \texttt{seed} field. This guarantees that a given seed always produces the same plant
arrangement and geometry, while different seeds produce meaningfully different visual results. The RNG is
also used in ecosystem placement, ensuring that the full scene is reproducible from a single seed.

% ──────────────────────────────────────────────────────────────────────────────────────────────────────────
\section{Turtle Graphics}
\label{sec:turtle}

\subsection{3D Turtle State}

The turtle maintains two unit vectors: \texttt{heading} ($\mathbf{H}$), the direction of travel, and
\texttt{normal} ($\mathbf{N}$), the up-vector of the current frame. The binormal $\mathbf{L} = \mathbf{H}
\times \mathbf{N}$ is computed on demand to complete the right-handed orthonormal frame. Position is tracked
as the most recent point in an internal path buffer. \texttt{Roll}, \texttt{Turn}, and \texttt{Pitch} actions
rotate the frame about $\mathbf{H}$, $\mathbf{L}$, and $\mathbf{N}$ respectively. \texttt{Push} saves the
full per-branch state (position, heading, normal and tracking values for environmental variation) to a
\texttt{StackEntry} on an internal stack. \texttt{Pop} restores the saved state.

\subsection{Branch Geometry}

A branch step (\texttt{Action::Branch(length, radius)}) appends a cylinder to the mesh. The cylinder is
generated by sweeping a circle of \texttt{lod\_segments} vertices along the current heading direction for
\texttt{length} units. Branch taper is applied by linearly scaling the radius from base to tip by the
\texttt{taper} parameter. UV coordinates are set to the sentinel value \texttt{[-1, -1]} to distinguish
branch geometry from leaf geometry in the fragment shader, where it suppresses the leaf oval mask.

\subsection{Leaf Geometry}

A leaf step (\texttt{Action::Leaf(width, height)}) appends a planar quad attached at the current turtle
position, oriented perpendicular to the current heading. The leaf texture is a $3 \times 3$ atlas of nine
distinct leaf variants. Each leaf quad samples a randomly selected cell from the atlas: a random integer $c
\in [0, 8]$ determines the column ($c \bmod 3$) and row ($\lfloor c / 3 \rfloor$), and the UV corners are
mapped to the corresponding $\frac{1}{3} \times \frac{1}{3}$ sub-region. Both the front and back faces use
the same atlas cell. This produces visual variety across leaves of the same species without additional
geometry or per-leaf parameters. In the fragment shader, leaf fragments are identified by a non-negative UV
$x$ coordinate; an opacity texture is sampled at the leaf UV, and fragments with opacity below 0.1 are
discarded, producing a per-pixel alpha cutout that trims the quad to the shape of each leaf variant.

\subsection{Cut Symbol}

The \texttt{Cut} action terminates the current branch by skipping all subsequent actions up to the matching
\texttt{]}~\cite{prusinkiewiczAlgorithmicBeautyPlants1990}. This is implemented in \texttt{do\_actions()} by
tracking bracket depth separately from the turtle stack: when \texttt{Cut} is encountered inside a
\texttt{[...]} block, a \texttt{cut\_at} depth counter is set, and all subsequent \texttt{Push}/\texttt{Pop}
adjust only the depth counter (not the stack) until the matching depth is reached and state is restored. This
avoids corrupting the turtle stack while correctly implementing the cut semantics.

\subsection{Space Pruning}

To prevent branches of one plant from occupying the same region of space as already existing plants, the
turtle implements a voxel-grid occupancy check. The scene volume is split into axis-aligned cells of a
configurable size (\texttt{occupancy\_cell\_size}); each cell is identified by an integer $(x, y, z)$ tuple.

\medskip
Two separate hash sets track occupancy: \texttt{occupancy\_grid} holds cells recorded by fully built plants
and is read (but never written) while the current plant is being built; \texttt{current\_plant\_marks} holds
cells marked during the current plant's build and is written but never read during that same build. This
ensures that intra-plant branches do not prune each other, only branches that would enter a region already
claimed by a \emph{prior} plant are cut. After each plant is built, \texttt{reset()} flushes
\texttt{current\_plant\_marks} into \texttt{occupancy\_grid} so that subsequent plants see the accumulated
footprint. At the start of each full scene rebuild, \texttt{clear\_occupancy()} empties both sets. When
marking and checking cells every cell within the branches' radius from the base to tip is evaluated.

\medskip
Both hash sets use an FNV-1a hasher implementation in place of Rust's default SipHash. The hash-flooding
resistance of SipHash is unnecessary for keys that are already-uniform integers under internal control. For
the \texttt{(i32, i32, i32)} integer keys the simpler XOR-multiply hash reduces lookup and insert cost
significantly~\cite{fowlerFNVNonCryptographicHash2026,aumassonSipHashFastShortInput2012}.

\subsection{Environmental Deflections}

Six per-step deflections are applied during branch traversal:

\begin{itemize}
\item \textbf{Phototropism}: at each \texttt{Branch} step, the turtle frame is rotated toward the configured
  light direction by $\Delta\theta = k_\text{photo} \cdot \sin\theta$, where $\theta$ is the current angle
  between the heading and the light direction and $k_\text{photo}$ is the user-configured strength.
\item \textbf{Gravitropism}: the same formula applied toward the $+Y$ axis with a separate strength
  $k_\text{grav}$, modelling negative gravitropism (upward shoot growth against gravity).
\item \textbf{Wind}: a perturbation in the configured wind direction, scaled by \texttt{wind\_strength} and
  jittered by \texttt{wind\_turbulence} using the scene RNG, is applied as a static heading offset baked at
  geometry build time.
\item \textbf{Proprioception}: a simplified approximation of the Bastien et al.\ autocorrelation
  model~\cite{bastienUnifiedModelShoot2015}. A curvature accumulator increments with each deflection and
  decays at rate $\gamma$ per step, modelling the plant's tendency to straighten after bending. The
  accumulator is reset to zero on each \texttt{Push}, so proprioception acts independently within each
  branch.
\item \textbf{Heading drift}: a correlated stochastic perturbation applied at each \texttt{Branch} step. The
  drift vector is incremented by noise scaled by \texttt{variation\_noise\_strength} and decayed by
  \texttt{variation\_decay\_rate}, producing the continuous meandering of real branch paths with memory
  across consecutive steps (in contrast to independent per-step jitter).
\item \textbf{Per-branch colour drift}: a similar drift process is applied to the RGB colour offset, producing
  natural colour variation across branches of the same plant without any grammar-level colour rules. Both
  drift processes are reset on \texttt{Push}.
\end{itemize}

% ──────────────────────────────────────────────────────────────────────────────────────────────────────────
\section{Plant Grammars}
\label{sec:grammars}

Eight species were implemented in total, descriptions of any unique symbols and parameters as well as figures
and full grammars are available in Appendix \ref{app:plants}.

\subsection{Fern}

The fern grammar builds a crown-based frond structure (Appendix~\ref{app:fern}). The crown symbol
$\mathrm{Cy}(e, \delta)$ waits $e$
steps then expands radially into three fronds at $120°$ intervals before spawning a deeper crown layer. Each
frond grows through $\mathrm{L}(a, d, \delta)$ nodes which branch once mature ($a \ge a_b$), provided there
is room to grow in both node distance and depth. The droop angle $\varphi$ applied at each branch node
increases with depth and node distance, creating the characteristic arching frond shape. Branch segments
$\mathrm{Br}(a, \delta)$ age in place, scaling the rendered geometry by age and depth. The growth and
structure of the fronds is based on Campbell's approach~\cite{campbellDescribingShapesFern1996}.

\subsection{Mint}

Mint has pairs of opposite leaves at each internode, with each successive pair rotated $90°$
(Appendix~\ref{app:mint}). The apex $A(t)$ increments a node counter and emits two internode leaf-pair units
each step. Every third node ($t \equiv 2 \pmod{3}$) two opposite sub-stems $B$ are inserted, each of which
repeats the leaf pattern. Growth terminates once the node count reaches $n_A$, after which further apex
rewrites omit the continuation.

\subsection{Wildflower}

The wildflower uses a stochastic branching grammar (Appendix~\ref{app:wildflower}). The plant axiom $P$ emits
a single apex $A(r)$ with a random starting jitter seed. Each step the apex extends three stem segments, spawns
a pair of lateral branches via $I(40)$, and continues with increased jitter $A(a{+}1)$. $I(a)$ immediately
resolves into two opposite branch tips $Q$ at ${\pm}a°$, where each $Q$ either extends stochastically ($p=0.5$)
or terminates by emitting a further $I(r)$ sub-branch. Flowers are appended to each apex step when in season.

\subsection{\textit{Lychnis}}

The \textit{Lychnis} grammar is based on the grammar described on Figure 3.14 in \textit{The Algorithmic
Beauty of Plants}~\cite{prusinkiewiczAlgorithmicBeautyPlants1990} (Appendix~\ref{app:lychnis}). It models a
campion with opposite leaves and a terminal flower. The apex $\mathrm{A}(d)$ waits $d_{\text{delay}}$ steps
before expanding, producing two opposite leaf pairs and two lateral apices at staggered delays to create
asynchronous branching. A terminal internode $\mathrm{I}(10)$ follows, emitting a flower $\mathrm{K}(0)$ only
when the seasonal parameter $\sigma$ falls within the bloom window. Leaves and flowers grow incrementally via
their size counters and are rendered scaled to $s_{\max}$.

\subsection{\textit{Mycelis}}

Adapted from L-System 3.3 in \textit{The Algorithmic Beauty of
Plants}~\cite{prusinkiewiczAlgorithmicBeautyPlants1990} (Appendix~\ref{app:mycelis}), the model for
\textit{Mycelis} is one of the more complex L-Systems. It makes heavy use of context-sensitive rules to send
signals throughout the plant to trigger developmental changes. Acropetal flowering is driven by a basipetal
$S$ signal that propagates through the main axis via $I(t)\!\to\!S$, triggering an apex $A$ to flower when
$S$ or $V$ appears to its left. Lateral buds $G$ are activated once per internode by the context signal $T$,
growing into a new internode--apex pair. The signal chain $S\!\to\!T\!\to\!W\!\to\!V$ advances one step each
generation, ensuring acropetal ordering. This progression is visible in Figure~\ref{fig:mycelis} where
branches and flowers propagate down the plant.

\subsection{\textit{Capsella}}

The \textit{Capsella} grammar is derived from Figure 3.5 in \textit{The Algorithmic Beauty of
Plants}~\cite{prusinkiewiczAlgorithmicBeautyPlants1990} (Appendix~\ref{app:capsella}). The vegetative apex
$a(t)$ counts down from $n_{\mathrm{veg}}$, producing leaves with golden-angle phyllotaxis on each step,
before transitioning to the flowering apex $A$.

\subsection{Tree}

The tree uses a ternary branching grammar with parametric internode growth and deciduous leaf shedding
(Appendix~\ref{app:tree}). Two apex symbols share the same branching template: the trunk apex $T$ extends
three internodes before the branch node, while the branch apex $X$ extends one. Both then emit three lateral
branches at roll offsets $-\beta$, $-4\beta$, and $-9\beta$, each leaned by $-\beta$ and followed by a leaf
cluster $G$. Leaf clusters render four differently oriented leaves; each leaf is independently dropped with
a season-dependent probability $p_{\mathrm{shed}}$ once the plant exceeds a minimum age. Bark colour
interpolates from a young to old value as the plant ages and leaf colour blends to amber-gold before shedding.

\subsection{Bush}

The bush uses the same three-branch template as the tree but with simpler geometry (Appendix~\ref{app:bush}):
all internodes are a fixed length (no parametric growth), leaf quads are placed inline at branch tips rather
than in a separate cluster symbol, and seasonal variation is handled by scaling leaf dimensions rather than
stochastic shedding. The trunk symbol $T$ extends one segment before delegating to the first-level brancher
$X$; the recursive sub-branch symbol $N$ extends two segments and interleaves leaf pairs at each sub-node,
producing a denser rounded crown.

% ──────────────────────────────────────────────────────────────────────────────────────────────────────────
\section{Environmental Response}

The per-branch deflections (phototropism, gravitropism, wind, proprioception, and OU drift) described in
Section~\ref{sec:turtle} are applied at the geometry level by the turtle interpreter. Higher-level
environmental response --- seasonal dormancy, leaf shedding, flowering phenology, and colour change --- is
encoded directly in the parametric production closures of each plant grammar, making it transparent to the
rewriting engine.

\subsection{Seasons}
\label{sec:seasons}

\begin{figure}[h]
\centering
\resizebox{\textwidth}{!}{\includegraphics{seasons}}
\caption{Mature forest scene throughout the year, clockwise from top left: spring, summer, autumn, winter.}
\label{fig:seasons}
\end{figure}

The season is a single \texttt{f32} value between $0$ and $1$: 0 corresponds to mid-winter, 0.25 to spring,
0.5 to summer, and 0.75 to autumn. Herbaceous perennial species implement a dormancy model in which the
effective L-System age is reduced by a dormancy factor $d(\mathit{season})$:

\begin{equation}
\mathit{age\_eff} = \bigl\lfloor \mathit{age} \cdot (1 - d) \bigr\rceil
\end{equation}

The dormancy factor $d$ is computed by a smooth-step curve that ramps from zero at summer to a
species-specific maximum at peak winter, plateaus through mid-winter, then recovers to zero by spring.
Scaling the age rather than individual rule parameters keeps the grammars independent of season.

\medskip
Autumn leaf colouration is implemented as a linear interpolation between the summer leaf colour and a
species-specific dormant colour. The interpolated colour is computed once inside \texttt{generate()} and
captured by the parametric closure, requiring no per-frame recomputation.

\medskip
Flowering phenology is implemented by gating the presence of flower symbols on a \texttt{show\_flower}
boolean computed from the current season. For species that retain seed heads after flowering
(e.g.\ wildflower), a separate \texttt{show\_seed} gate models the phenological transition from flower to
seed with a reduced symbol size.

\medskip
The global season tint, interpolated from white at summer through warm amber in autumn to cool grey in
winter, is written to the uniform buffer on every frame via \texttt{season\_tint()} without triggering a
geometry rebuild, enabling smooth real-time colour blending as the season slider is adjusted.

% ─────────────────────────────────────────────────────────────────────────────
\section{Ecosystem Placement}
\label{sec:ecosystem}

\subsection{Kernel-Based Placement}

Plant positions are generated by sampling from a 2D probability density function (PDF) defined over a $32
\times 32$ grid spanning the scene area. The grid is initialized with uniform probability. Each time a plant
is placed, a deformation kernel centred on that plant's position updates the PDF
~\cite{laneGeneratingSpatialDistributions2002}:

\begin{itemize}
\item \textbf{Inhibitory}: a Gaussian bell curve reduces the probability near the placed plant, modelling
  competitive exclusion.
\item \textbf{Promotional}: the bell curve increases probability, modelling facilitation (e.g.\ nurse plants).
\item \textbf{Mixed}: a negative bell at close range combined with a weak positive lobe at medium range.
\item \textbf{Neutral}: no kernel applied; produces a uniform random distribution as a baseline.
\end{itemize}

Species are sampled from a user-configured weighted list. The weight determines relative abundance: equal
weights produce a mixed community, skewed weights produce a dominant species with subordinates.

\subsection{Self-Thinning}

After placement, optional self-thinning iteratively removes the smaller plant from any pair closer than the
configured thinning radius. The algorithm continues until no overlapping pairs remain, producing an
overdispersed (regular) distribution characteristic of mature single-species
stands~\cite{laneGeneratingSpatialDistributions2002}.

\subsection{Succession}

The optional succession algorithm alternates between growth steps, increasing each plant's radius increases
by its species' intrinsic growth rate, and competition steps, removing the plant with lower shade tolerance
from overlapping pairs. Each surviving plant has a 25\% chance of producing an offspring nearby, inheriting
the parent's species. This models the succession of pioneer communities (fast-growing, low shade tolerance)
toward climax communities (slow-growing, high shade tolerance)~\cite{laneGeneratingSpatialDistributions2002}.
Shade tolerance and growth rate are intrinsic properties of each \texttt{PlantType} rather than
user-adjustable parameters, reflecting their biological basis.

% ──────────────────────────────────────────────────────────────────────────────────────────────────────────
\section{Level of Detail}
\label{sec:lod}

Each plant is assigned an LOD tier based on its distance from the camera:

\begin{table}[h]
\centering
\begin{tabular}{|l|l|l|l|}
  \hline
  \textbf{Tier}   & \textbf{Distance} & \textbf{Rendering} \\
  \hline
  \texttt{Near}   & $<80$ m           & 8 cylinder segments, all branches, all leaves \\
  \texttt{Mid}    & $80$--$160$ m     & 5 cylinder segments, all branches, reduced leaf density \\
  \texttt{Far}    & $160$--$300$ m    & 3 cylinder segments, $r \geq 0.005$ only, reduced leaf density \\
  \texttt{Beyond} & $>300$ m          & Mesh skipped; line skeleton only, no leaves \\
  \texttt{Culled} & Outside frustum   & Nothing rendered \\
  \hline
\end{tabular}
\caption{LOD tier thresholds (default values; user-adjustable).}
\label{tab:lod}
\end{table}

\subsection{Leaf-Skip Probability}

Each LOD tier carries a configurable \emph{leaf-skip probability}: \texttt{leaf\_skip\_near},
\texttt{leaf\_skip\_mid}, and \texttt{leaf\_skip\_far} in \texttt{LodSettings}. At each
\texttt{Action::Leaf} step the turtle draws a uniform random number and skips the quad if it is below
the skip probability for the tier containing it. \texttt{Beyond} and \texttt{Culled} plants always skip
their leaves. Using a fixed limit and skipping leaves once the limit is reached results in an obvious and
unnatural pattern with some regions being bare whilst their neighbours remain untouched. Additionally,
escalating the probability with the LOD tiers results in the most aggressive culling taking place in the
least visible areas.

\subsection{Frustum Culling}

Plants entirely outside the camera view frustum are assigned the \texttt{Culled} tier and contribute no
geometry. The view frustum is extracted from the view-projection matrix using the Gribb--Hartmann
method~\cite{gribbFastExtractionViewing}, yielding six half-space planes. Each plant is then tested as a
bounding sphere of configurable radius \texttt{cull\_radius}. Plants are only culled if their sphere lies
fully outside at least one plane.

\subsection{Adaptive LOD Pressure}

A \texttt{lod\_pressure} multiplier between $0.1$ and $1.0$ is used to scale the LOD distance thresholds and
increase the leaf-skip probabilities. When the assembled scene geometry exceeds the GPU buffer limit the
multiplier is reduced causing the LOD boundaries to shrink and the leaf-skip probabilities to grow,
$p \leftarrow \max(0.1,\; p \times 0.75)$. Once the buffer fill ratio falls bellow 65\% the pressure multiplier
is gradually restored: $p \leftarrow \min(1.0,\; p / 0.9)$.

A change in \texttt{lod\_pressure} triggers a full rebuild so that the tightened thresholds take effect
on the next frame. Recovery is slower than degradation and requires significant headroom in order to prevent
oscillation between overflow and fit states. The current pressure value is displayed in the debug HUD when
it falls below 1.0.

% ──────────────────────────────────────────────────────────────────────────────────────────────────────────
\section{Rendering Pipeline}

\subsection{Geometry Buffers}

The mesh pipeline uses four per-vertex buffers: position (\texttt{Vec3}), colour (\texttt{Vec3}), normal
(\texttt{Vec3}), and UV (\texttt{[f32; 2]}). All plants in a scene are combined into a single set of buffers
before upload. Line geometry is stored in a separate buffer pair. The uniform block is 96 bytes: 64 bytes
for the 4$\times$4 view-projection matrix, 12 bytes for the light direction, 4 bytes for the ambient factor,
12 bytes for the season tint colour, and 4 bytes of padding to satisfy WGSL alignment requirements.

\subsection{Shaders}
\label{sec:shaders}

The \textbf{mesh shader} (\texttt{mesh.wgsl}) implements Lambert diffuse lighting with tangent-space normal
mapping for leaves. In the vertex stage the clip position is computed by applying the view-projection matrix
and colour, world normal, UV, and tangent are passed through as interpolants. In the fragment stage:

\begin{enumerate}
\item The UV sentinel ($\mathit{uv}.x < 0$) distinguishes branch faces from leaf faces.
\item For leaf fragments, a colour texture, normal map, and opacity texture are sampled from the leaf atlas.
  Fragments with opacity below 0.1 are discarded, producing per-pixel cutout silhouettes.
\item The tangent-space normal is decoded from $[0,1]^3$ to $[-1,1]^3$ and transformed to world space using a
  Gram--Schmidt re-orthogonalised tangent frame ($T' = \text{normalise}(T - (T \cdot N)N)$, $B = N \times T'$).
\item Lambert diffuse is computed as $\max(\mathbf{n} \cdot \mathbf{L}_\mathit{dir},\ 0)$ and combined with
  the ambient factor: $\mathit{lighting} = \mathit{ambient} + (1 - \mathit{ambient}) \cdot \mathit{diffuse}$.
\item The final colour is $\mathit{base\_colour} \times \mathit{season\_tint} \times \mathit{lighting}$,
  where $\mathit{base\_colour}$ is the texture colour modulated by vertex colour for leaves, or the vertex
  colour alone for branches.
\end{enumerate}

The \textbf{line shader} (\texttt{lines.wgsl}) is a minimal pass-through: it transforms position by the
view-projection matrix and outputs the interpolated vertex colour unchanged. The season tint is not applied
to line geometry.

% ──────────────────────────────────────────────────────────────────────────────────────────────────────────
\section{Testing and Validation}
\label{sec:testing}

\subsection{Unit Tests}

The project includes a suite of unit tests across five modules.

\subsubsection{L-System interpreter}
The classical D0L algae grammar ($A \to AB$, $B \to A$) is verified to produce the correct Fibonacci string
sequence across the first eight generations. Terminal symbol pass-through is confirmed, as is the distinction
between symbols without rules (which copy unchanged) and those with rules (which rewrite). Stochastic rules
are tested at probabilities 0 and 1, and at probability 0 with a Normal fallback rule to verify the fallback
chain. Context-sensitive matching is covered in depth: left-only and right-only contexts, failure at the
string boundary, failure when the neighbour has the wrong identity, correct skipping of
\texttt{Ignore}-marked symbols, correct skipping of bracketed sub-expressions during left-context traversal,
right-context search across multiple branches. Parametric rules are tested for reading inline symbol values
and for matching by discriminant regardless of parameter value; an \texttt{Ignore}-typed parametric symbol is
confirmed not to satisfy a context requirement. Priority ordering is verified: context-sensitive rules fire
before parametric rules and before stochastic rules for the same symbol, and a normal fallback fires when
context goes unmatched.

\subsubsection{Turtle graphics}
Geometry generation is tested for both branch and leaf primitives: a \texttt{Branch} action with 8 segments
produces 16 vertices and 48 indices; with 4 segments it produces 8 vertices and 24 indices; with 0 segments
or a radius below the minimum-radius threshold no mesh is generated. A \texttt{Leaf} action produces a
double-sided quad (8 vertices, 12 indices), and multiple leaves accumulate correctly.

\medskip
Orientation is verified for all three rotation axes: \texttt{Turn} rotates the heading around the normal,
\texttt{Roll} changes only the normal (leaving the heading unaffected), and \texttt{Pitch} tilts both vectors
around the binormal. Scale is confirmed to multiply travel distance proportionally.

\medskip
State management tests verify that a single \texttt{Push}/\texttt{Pop} restores position and heading after
branching, and that depth-two nesting fully restores the outer frame on both pops. The \texttt{Cut} symbol is
tested inside a branch block (geometry is skipped), at depth zero (ignored), and spanning only the directly
enclosing bracket (geometry after the closing bracket still executes). The leaf-skip probability is confirmed
to suppress all leaf geometry when set to 1.0.

\medskip
Space pruning tests cover: intra-plant branches never blocking each other (marks are in a separate set not
checked during the same build); inter-plant blocking after \texttt{reset()} flushes marks into the shared
grid; mid-segment overlap detection (a branch whose \emph{tip} lands in a clear cell is still cut if its path
crosses an occupied cell); radius-based volumetric expansion blocking a cylinder that overlaps the marked
region without sharing the centre-line cell; and well-separated branches not being pruned. The
disabled-pruning case confirms all branches accumulate when the feature is off.

\medskip
Geometry combination tests confirm that \texttt{combine\_line\_geometries} and
\texttt{combine\_mesh\_geometries} correctly offset indices across multiple geometries and that
\texttt{mesh\_index\_count} matches the accumulated index buffer length.

\subsubsection{Season utilities}
Tests verify \texttt{season\_name} across all four quadrants (including wrap-around and negative inputs);
\texttt{season\_tint} returning white at summer and the configured winter tint at winter, with continuity
across the summer boundary; and \texttt{season\_needs\_rebuild} returning true across a threshold boundary
and false within the same bucket. Generation counter tests confirm that \texttt{mark\_dirty()} increments the
counter and that initial equality with \texttt{u64::MAX} forces the first-frame rebuild. \texttt{PlantType}
property tests verify that \texttt{shade\_tolerance} is in \texttt{[0, 1]} and \texttt{growth\_rate} is
positive for all nine species.

\subsubsection{Ecosystem placement}
Kernel helper functions are tested at boundary values and for monotonicity. The \texttt{mixed\_kappa} kernel
is verified to be negative at the origin (inhibitory inner lobe), to switch sign at the inner-to-outer
boundary ($r = 0.4$), and to decay to near zero at the outer edge. The generator is tested for: empty species
list returning no plants; zero plant count returning no plants; correct plant count when species and count
are given; seed reproducibility (same seed produces identical positions and types); different seeds producing
different arrangements; all generated positions lying within the configured area bounds; species weights
being respected; and all plant instances being placed at $y = 0$. The \texttt{self\_thinning} function is
unit-tested directly for edge cases: an empty input remaining empty; a single plant being unchanged; two
non-overlapping plants both surviving; and an overlapping pair retaining only the plant with the larger radius.

\subsubsection{Plant trait and utilities}
The \texttt{dormancy\_factor} function is tested across the full seasonal cycle: zero dormancy before onset,
the smooth-step ramp at the onset and recovery midpoints, the plateau at peak winter, wrap-around at season
1.0, and the effect of a positive phase offset shifting onset earlier. The \texttt{lerp\_bark\_colour}
function is tested at age zero, at \texttt{max\_age}, at the midpoint, and clamped beyond \texttt{max\_age}.

\medskip
\texttt{PlantInstance} construction is tested for correct field assignment via \texttt{with\_transform}, and
preservation of \texttt{delay\_years} across transform calls. Cloning is confirmed to copy all fields
including position, scale, rotation, delay, and plant type.

\medskip
Smoke tests confirm that all nine plant types produce non-empty action sequences at iteration 0 and iteration
1. Additional tests verify that \texttt{set\_iteration} with the same value does not change the action count
and that \texttt{years\_per\_iteration} and \texttt{shade\_tolerance} are within their valid ranges.

\subsection{Visual Validation}

\begin{itemize}
\item \textit{Mycelis}: the flowering state propagates acropetally (from tip to base) as age increases
  (Figure~\ref{fig:mycelis}).
\item \textit{Capsella}: lateral branches appear at every other node with alternating roll offsets
  (Figure~\ref{fig:capsella}).
\item Tree: no two instances of the same tree grammar at the same age are geometrically identical due to
  per-branch roll jitter, confirming that the stochastic variation is working as intended.
\item Cut symbol: at sufficient age, a fraction of leaf clusters in the tree are absent, and the associated
  branches terminate at the push point rather than continuing to the leaf.
\item Seeded reproducibility: in addition to the unit tests, changing and restoring the seed was visually
  confirmed to restore the previous state.
\end{itemize}
